\section*{Bab \ref{ch:b3:pericyclic} --- Reaksi Perisiklik}

\begin{solution}{exo:b3:pericyclic:1}
Diels--Alder: \omterm{def:b3:pericyclic:families}{sikloadisi} [4+2]. Pembukaan siklobutena: elektrosiklik. Cope: sigmatropik
[3,3]. Sulfolena yang melepaskan $\ce{SO2}$: keletropik. Ozon dan alkena: \omterm{def:b3:pericyclic:dipolar}{sikloadisi dipolar-1,3}
([3+2]). Pergeseran $[1,5]$-H: sigmatropik.
\end{solution}

\begin{solution}{exo:b3:pericyclic:2}
Heksatriena (6 elektron): termal \omterm{def:b3:pericyclic:rotation-modes}{disrotatori}, fotokimia \omterm{def:b3:pericyclic:rotation-modes}{konrotatori}.
Butadiena (4): termal \omterm{def:b3:pericyclic:rotation-modes}{konrotatori}, fotokimia \omterm{def:b3:pericyclic:rotation-modes}{disrotatori}.
\end{solution}

\begin{solution}{exo:b3:pericyclic:3}
Empat elektron, termal: pembukaan \omterm{def:b3:pericyclic:rotation-modes}{konrotatori}. Kedua gugus metil, yang berada pada muka
cincin yang berlawanan, sama-sama berputar ke luar: (\textit{E},\textit{E})-heksa-2,4-diena.
\end{solution}

\begin{solution}{exo:b3:pericyclic:4}
Reaksi [2+2] \omterm{def:b3:pericyclic:faciality}{suprafasial} memiliki $4q$ elektron: dalam keadaan dasar, HOMO satu
etena dan LUMO etena yang lain bertumpang tindih dengan satu ujung ikatan dan satu ujung antiikatan.
Dalam keadaan tereksitasi, $\pi^*$ yang terisi tunggal pada satu molekul berperan sebagai
HOMO, dan kedua ujungnya cocok.
\end{solution}

\begin{solution}{exo:b3:pericyclic:5}
\omterm{def:b3:point-groups:mirror}{Bidang cermin} dipertahankan. Heksatriena: $\psi_1$ S, $\psi_2$ A, $\psi_3$ S terisi; $\psi_4$ A, $\psi_5$ S, $\psi_6$ A kosong.
Sikloheksadiena, menurut urutannya: $\sigma$ S, $\pi_1$ S, $\pi_2$ A terisi; $\pi_3^*$ S, $\pi_4^*$ A, $\sigma^*$ A kosong. Dengan menghubungkan menurut
simetri secara berurutan, $\psi_1 \to \sigma$, $\psi_2 \to \pi_2$, $\psi_3 \to \pi_1$: terisi menuju terisi, penutupan \omterm{def:b3:pericyclic:rotation-modes}{disrotatori}
diizinkan secara termal.
\end{solution}

\begin{solution}{exo:b3:pericyclic:6}
Enam elektron di bawah cahaya: \omterm{def:b3:pericyclic:rotation-modes}{konrotatori}, sehingga kedua gugus metil ujung berakhir pada
muka yang berlawanan: \textit{trans}-5,6-dimetilsikloheksa-1,3-diena (reaksi termal memberikan
isomer \textit{cis}).
\end{solution}

\begin{solution}{exo:b3:pericyclic:7}
Hidrogen pada C5 berpindah ke C1 diena, tetangganya di sekeliling cincin,
melalui keadaan transisi enam elektron ($\sigma2_s + \pi4_s$), secara \omterm{def:b3:pericyclic:faciality}{suprafasial}: satu migrasi
memberikan 1-metil-, migrasi kedua 2-metilsiklopentadiena. Rantai lima karbon memungkinkan
hidrogen tetap berada pada satu muka, yang memang dipaksakan oleh cincin yang kecil itu.
\end{solution}

\begin{solution}{exo:b3:pericyclic:8}
$\pi2_s$ (alkena) $+ \pi2_a$ (C=C keten, yang diserang pada muka yang berlawanan melalui
$\pi^*$ C=O yang tegak lurus): jumlah komponen $(4q + 2)_s$ adalah satu (alkena), ganjil:
diizinkan secara termal.
\end{solution}

\begin{solution}{exo:b3:pericyclic:9}
Dalam kursi itu terdapat C=C gugus vinil (atom 1--2), oksigen (3), karbon
$\ce{OCH2}$ (4), dan C=C alil (5--6); ikatan O--C putus dan ikatan C1--C6 terbentuk:
produknya adalah pent-4-enal,
$\ce{OHC-CH2-CH2-CH=CH2}$.
\end{solution}

\begin{solution}{exo:b3:pericyclic:10}
Delapan elektron adalah $4q$; dengan satu komponen \omterm{def:b3:pericyclic:faciality}{antarafasial}, susunannya bertopologi
Möbius, yang aromatik dengan $4q$ elektron: keadaan transisinya aromatik dan
reaksinya diizinkan secara termal (seperti penutupan \omterm{def:b3:pericyclic:rotation-modes}{konrotatori} oktatetraena).
\end{solution}

\begin{solution}{exo:b3:pericyclic:11}
Memasangkan koefisien-koefisien terbesar menghubungkan N3 dengan karbon ujung: triazola
1,4-tersubstitusi. Tanpa katalis, koefisien-koefisiennya hanya sedikit berbeda, sehingga
isomer 1,4 dan 1,5 sama-sama terbentuk; tembaga(I) membuat reaksinya bertahap melalui
tembaga asetilida dan hanya memberikan isomer 1,4.
\end{solution}

\begin{solution}{exo:b3:pericyclic:12}
Hidrogen menjembatani ujung-ujung SOMO radikal $j$ atom, $k = (j + 1)/2$. Koefisien ujungnya
sebanding dengan $\sin(k\pi/(j + 1))$ dan $(-1)^{k+1}\sin(k\pi/(j + 1))$. Untuk $j = 7$, $k = 4$: tanda berlawanan, sehingga
cuping-cuping yang bertanda sama berada pada muka yang berlawanan di kedua ujung: \omterm{def:b3:pericyclic:faciality}{antarafasial}. Untuk
$j = 5$, $k = 3$: tanda sama pada muka yang sama: \omterm{def:b3:pericyclic:faciality}{suprafasial}.
\end{solution}

\begin{solution}{pb:b3:pericyclic:1}
\textbf{1.} Ikatan $\sigma$ C9--C10 cincin B; bersama 5,7-diena, ikatan itu memberikan heksatriena,
C10=C5--C6=C7--C8=C9.

\textbf{2.} Enam elektron (kedua ikatan $\pi$ dan ikatan $\sigma$): pembukaan cincin elektrosiklik.

\textbf{3.} Secara termal, sikloheksadiena dan heksatriena terbukanya hanya mencapai kesetimbangan
dengan lambat, melalui penghalang yang tinggi, dan cincin itulah sisi yang lebih stabil;
diena itu menyerap sinar ultraviolet, dan keadaan tereksitasinya membuka dalam
orde pikodetik.

\textbf{4.} \omterm{def:b3:pericyclic:rotation-modes}{Konrotatori}: dalam keadaan tereksitasi, LUMO sistem triena yang terisi tunggal
($k = 4$) memiliki ujung-ujung yang berlawanan tanda.

\textbf{5.} Ikatan C6=C7 berasal dari cincin: kedua lanjutan rantainya, C5 dan C8,
berada pada sisi yang sama terhadap ikatan itu, dan tetap demikian: \textit{Z}.

\textbf{6.} Pembukaan termal, yang hanya diizinkan secara \omterm{def:b3:pericyclic:rotation-modes}{disrotatori}, memiliki penghalang yang jauh melampaui
apa yang dapat disediakan oleh panas tubuh, dan akan menanjak ke arah triena yang kurang stabil.

\textbf{7.} Dari C19 (metil pada C10) ke C9: pergeseran sigmatropik $[1,7]$-H.

\textbf{8.} Delapan: keenam elektron $\pi$ triena dan kedua elektron ikatan C--H.

\textbf{9.} $\sigma2 + \pi6$. Secara \omterm{def:b3:pericyclic:faciality}{suprafasial} keduanya akan bertipe s, dengan satu komponen $(4q + 2)_s$ dari masing-masing:
genap, terlarang. Dengan triena \omterm{def:b3:pericyclic:faciality}{antarafasial}, $\sigma2_s + \pi6_a$: satu komponen yang dihitung, ganjil,
diizinkan.

\textbf{10.} Tujuh atom dalam rantai berbentuk heliks yang berkonfigurasi \textit{Z} memungkinkan hidrogen berpindah dari satu muka ke
muka yang lain; rantai tiga atom tidak dapat terpuntir sejauh itu.

\textbf{11.} Möbius (satu komponen \omterm{def:b3:pericyclic:faciality}{antarafasial}) dengan delapan elektron, $4q$: aromatik.

\textbf{12.} Pergeseran itu diizinkan secara termal; penghalangnya dilewati dengan lambat pada suhu tubuh.

\textbf{13.} Enam elektron di bawah cahaya: \omterm{def:b3:pericyclic:rotation-modes}{konrotatori} lagi, tetapi penutupan itu dapat memutar ujung-ujungnya ke
arah yang lain dan menempatkan H9 dan metil C19 pada muka-muka yang lain: sebuah stereoisomer
cincin tertutup, yaitu lumisterol.

\textbf{14.} Dengan ikatan tengah \textit{E}, C19 dan C9 berada pada sisi rantai yang berlawanan dan saling
berjauhan: hidrogen tidak dapat menjangkaunya.

\textbf{15.} Pravitamin D$_3$ sendiri menyerap cahaya dan diubah menjadi lumisterol dan takisterol,
yang tidak memberikan vitamin D; tercapailah campuran fotostasioner, yang membatasi
produksinya.

\textbf{16.} Pembukaan terjadi dalam keadaan tereksitasi, dalam orde femtodetik sampai pikodetik;
pergeseran itu adalah reaksi termal dengan penghalang \qty{85}{kJ/mol}.

\textbf{17.} $t_{1/2} = \ln 2/k = 0.693/\qty{0.023}{h^{-1}} = \qty{30}{h}$.

\textbf{18.} $1 - \eu^{-0.023 \times 24} = 0.42$.

\textbf{19.} $\ln 10/k = 2.303/\qty{0.023}{h^{-1}} = \qty{100}{h}$.

\textbf{20.} $k(\qty{293.15}{K}) = k(\qty{310.15}{K})\,\eu^{-(E_a/R)(1/293.15 - 1/310.15)} = 0.023 \times \eu^{-1.91} = \qty{3.4e-3}{h^{-1}}$.

\textbf{21.} $2.303/\qty{3.4e-3}{h^{-1}} = \qty{680}{h}$, sekitar 28 hari.

\textbf{22.} Tahap termal itu tujuh kali lebih cepat pada \qty{37}{\celsius} daripada pada \qty{20}{\celsius}, sehingga
pravitamin yang dibuat pada suatu pagi di bawah sinar matahari terkonversi dalam beberapa hari.

\textbf{23.} Kerangka steroid, dengan banyak pusat stereonya, sudah tersedia dari
sterol itu; cahaya dan kehangatan lalu menjalankan kedua tahap perisiklik tepat seperti di dalam
kulit.

\textbf{24.} Pada suhu tubuh, 90\,\% pravitamin D$_3$ menjadi vitamin D$_3$ dalam sekitar
\qty{100}{h}, empat hari.
\end{solution}
